\section{An obstruction to reusing the original word enumerator}
\label{sec:altered-scheduler}

This observation concerns the proposed reinterpretation of the rejected
195-state table. Assume its changed positive fetch has the reported
source effect: with $L>0$ it decodes the seed $N+L-1$ and increments
$N$ by two. Retain the original fresh fetch: when $L=0$, decode seed
$N$ and increment $N$ by one. Write
\[
 p(a,b)=\frac{(a+b)(a+b+1)}2+b.
\]
The first unpaired component selects the command: zero is Shift, one is
modus ponens, and two is the first axiom command. The second component
is the residual $L$.

\paragraph{Claim.}
No command sequence beginning at a fresh fetch can begin with two
successive Shift commands.

\paragraph{Proof.}
Let $n$ be the fresh seed. If the first command is Shift, then
$n=p(0,b)$ for some $b\geq0$. Immediately afterwards $N=n+1$ and $L=b$.
If $b>0$, the changed positive fetch decodes
\[
 (n+1)+b-1=n+b
 =\frac{b(b+1)}2+2b
 =p(2,b-1).
\]
Its command selector is two, so its command is the first axiom command.
If $b=0$, then $n=0$ and the next fresh seed is one; since
$1=p(1,0)$, the next command is modus ponens. Neither case gives Shift.
\hfill$\square$

The original finite-word encoding theorem promises every requested
command prefix at a fresh candidate boundary. The claim gives a
length-two counterexample to that theorem for the altered scheduler.
Padding the terminal residual cannot repair this particular prefix,
because the calculation holds for every possible $b$.

This is not a proof that no different completeness argument could ever
use the changed scheduler. Such an argument would need a new compiler
that accommodates its restricted command streams, or an independent
simulation embedded away from a fresh boundary. It would also have to
establish sound initialization of the theorem store and a complete
physical source binding. None of those additional obligations follows
from observing a long nonhalting run of the rejected table.
